\section{Characterization details}
\label{sec:appx-char}

\subsection{Dense and expert-parallel MoE layers}
\label{sec:appx-arch}

Figure~\ref{fig:appx-arch} contrasts a dense Transformer layer with an
\ac{MoE} layer on two GPUs. Under \ac{TP}, a dense layer shards
attention heads and feed-forward weights evenly, so every GPU performs
the same work on all tokens and reaches each all-reduce at the same
time. In the \ac{MoE} layer, attention is still sharded by \ac{TP}, and
every GPU computes the same routing decision, but \ac{EP} places whole
experts on different GPUs. Each GPU then runs only the tokens routed to
its local experts. When the router sends more tokens to one GPU's
experts, the other GPU finishes early and waits at the all-reduce that
combines the expert outputs: this wait is the bubble of
Figure~\ref{fig:anatomy}.

\begin{figure}[t]
  \centering
  \inputplot{arch_bubble}
  \caption{\textbf{Where expert-parallel bubbles come from.} Time runs
    downward within each GPU. (a)~A dense layer under \ac{TP} splits
    work evenly, so no GPU waits. (b)~An \ac{MoE} layer under
    \ac{TP}$=$\ac{EP} places different experts on different GPUs; the
    router sends GPU~0's experts fewer tokens (squares), so GPU~0
    finishes early and idles until GPU~1 reaches the all-reduce that
    combines the expert outputs. Schematic; sizes are illustrative.}
  \label{fig:appx-arch}
\end{figure}

\subsection{Bubble size across models}
\label{sec:appx-context}

Figure~\ref{fig:context} shows how the per-layer bubble grows with
prefill size for the three \ac{MoE} models of \S\ref{sec:char-context},
in the default setup of \S\ref{sec:characterization}. From 8k to 128k
tokens, the median bubble grows from 0.22 to 6.41\,ms/layer on
Qwen3-30B-A3B, from 0.28 to 4.34\,ms/layer on Qwen1.5-MoE-A2.7B, and
from 0.20 to 2.70\,ms/layer on DeepSeek-V2-Lite.

\begin{figure}[t]
  \centering
   \inputplot{context_bubble}
  \caption{
    \textbf{Per-layer bubble duration grows with prefill size.}
    Median bubble across three \ac{MoE} architectures under BF16
    execution without load balancing. Error bars show min/max across
    repetitions.
  }
  \label{fig:context}
\end{figure}

\subsection{Lower-precision execution}
\label{sec:appx-fp8}

Lower-precision execution attacks a different part of the problem.
FP8 reduces the cost of expert computation, but does not change how
the router distributes expert work across GPUs. Nor, in our measured
regime, does it reduce expert computation enough to remove the
compute saturation during prefill: additional routed work on one GPU
still extends its completion time. Figure~\ref{fig:appx-fp8} compares
BF16 and FP8 in the load-balancing setup of Figure~\ref{fig:lb-sweep}
(DeepSeek-V2-Lite, SGLang, 8$\times$H100, \ac{EP}$=8$). At 128k,
moving from BF16 to FP8 changes the bubble from 4.79 to 4.81\,ms/layer
without load balancing and from 3.53 to 3.50\,ms/layer with LPLB, and
it changes by at most 0.05\,ms/layer at every other context. Faster
expert execution therefore does not, by itself, remove the GPU-to-GPU
timing differences caused by uneven routed work.

\begin{figure}[t]
  \centering
  \inputplot{precision_sweep}
  \caption{\textbf{Lower precision leaves the bubble unchanged.}
    Median per-layer bubble under BF16 and FP8 (hatched), without load
    balancing (Base) and with LPLB, in the setup of
    Figure~\ref{fig:lb-sweep}: DeepSeek-V2-Lite, SGLang, 8$\times$H100,
    \ac{EP}$=8$. Context denotes aggregate tokens across the 8-way
    \ac{EP} group.}
  \label{fig:appx-fp8}
\end{figure}
\subsection{Bubble predictability}
\label{sec:appx-predict}

We measure how well the router's token assignment predicts each GPU's
bubble on Qwen3-30B-A3B, served by \vllm on 4$\times$A100 with
\ac{TP}$=$\ac{EP}$=4$ and no expert load balancing, at five loads
(4k tokens at 8 requests/s, 8k at 4 and 2, 12k at 2, and 2k at 1;
400 \textsc{arXiv} requests each). For every prefill layer and GPU
(\num{193920} samples), we compare the measured bubble with two
predictors: the deployed router-aware predictor
(\S\ref{sec:design-predict}), and a static profile that assigns every
GPU the mean bubble of that layer from the first half of the run. Both
are evaluated on the second half. \Cref{fig:appx-predict} shows that
the router-aware predictor reduces the median absolute error from
\SI{123}{\micro\second} to \SI{8.9}{\micro\second} and the p95 from
\SI{452}{\micro\second} to \SI{45}{\micro\second}. With \ac{EPLB}
enabled, its median error is \SI{9.5}{\micro\second} (p95
\SI{43}{\micro\second}). Our probe resolves durations to
\SIrange{4}{8}{\micro\second}, so the router-aware median error is at
the measurement resolution.

\begin{figure}[t]
  \centering
  \inputplot{bubble_predict_error}
  \caption{\textbf{The router's assignment predicts each GPU's bubble.}
    Absolute error of the router-aware predictor and of a static
    per-layer profile on Qwen3-30B-A3B, 4$\times$A100,
    \ac{TP}$=$\ac{EP}$=4$, pooled over five serving loads.}
  \label{fig:appx-predict}
\end{figure}